\section{Introduction}
\label{sec:einleitung}

Large Language Models (LLMs) are increasingly integrated into applications that
read emails, summarise documents, retrieve web content, and execute actions on
behalf of users. This increases their usefulness but also exposes them to
content controlled by parties other than the application operator.

Prompt injection exploits this exposure by placing malicious instructions in
content processed by the model. If the model follows the injected instruction
instead of the trusted task, it may manipulate a response, disclose protected
information, or influence a connected tool
\parencites{perez2022ignore}{greshake2023indirect}. The 2025 OWASP Top 10 for
LLM Applications assigns prompt injection the identifier \emph{LLM01}, marking
it as the first category listed in the document
\parencite{owasp2025llm}.

Applications can encode developer instructions, user input, and tool results
using different roles or delimiters. These distinctions influence model
behaviour, but they do not constitute access control. Unlike a prepared SQL
statement, a prompt format does not give natural-language text a reliably
enforced type or authority. Security cannot therefore rely on formatting alone
to preserve the distinction between trusted instructions and untrusted data
\parencites{liu2023promptinjection}{wallace2024instruction}.

StruQ and SecAlign address this problem by training models to maintain an
instruction--data separation. Their original evaluations report substantial
reductions in attack success, particularly against manually constructed attacks
\parencites{chen2025struq}{chen2025secalign}. However, reported robustness is
difficult to compare across studies because evaluations differ in models,
tasks, attacker access, optimisation methods, and success criteria. Later
defence-adaptive studies further question whether results obtained under narrow
attack suites generalise to stronger attacker models
\parencites{jia2025critical}
             {yang2025checkpointgcg}
             {pandya2025astra}.

This creates two related questions: how much security can learned
instruction--data separation provide under the evaluated conditions, and how
should applications limit the consequences when the model still follows an
injected instruction?

\medskip
\noindent\textbf{Research Question} \quad
\textit{How effectively do learned instruction--data separation approaches
mitigate prompt injection under the evaluated attacker models, and what
system-level controls are required when they do not provide an independently
enforced security boundary?}

\medskip
To answer this question, the paper compares the original StruQ and SecAlign
evaluations with later defence-adaptive studies and relates their limitations
to system-level security controls. No original experiments or independent
replications are conducted.

Chapter~\ref{sec:llms} introduces LLM prompting and structured context.
Chapter~\ref{sec:prompt-injection} describes prompt-injection delivery and
construction methods, and Chapter~\ref{sec:threat} defines the threat model used
for the comparison. Chapter~\ref{sec:struq} examines StruQ and SecAlign, while
Chapter~\ref{sec:limits} evaluates the limits of their reported robustness.
Chapter~7 derives practical deployment implications, and
Chapter~\ref{sec:conclusion} answers the research question.